% !PRINTED 263
\weberpage{263}
Wir kommen also zu dem ersten Resultate:

1.~\gesperrt{Eine Substitution $A$ ist immer auf die Normalform $S$ transformirbar, wenn $a+d$ nicht $=\pm 2$
ist. $\sigma$ ist reell oder imaginär, je nachdem $(a+d)^2 - 4$ quadratischer Rest oder Nichtrest
von $p$ ist.}

In dem noch übrigen Falle, wo $a+d = \pm 2$ ist, lässt sich
die Normalform $S$ nicht herstellen; dagegen können wir $A$ in
diesem Falle auf eine andere Normalform, nämlich
\weq{6}{T = \begin{pmatrix} 1, & t \\ 0, & 1 \end{pmatrix}}
bringen, in der $t$ reell ist.

Um dies zu beweisen, können wir zunächst
\weq{7}{a + d = +2}
annehmen, weil der andere Fall $a+d = -2$ durch Aenderung
aller Vorzeichen von $a, b, c, d$ auf diesen zurückgeführt wird.

Aus (7) aber folgt noch
\weq{8}{a - 1 = -(d-1), \quad (a-1)(d-1) = b\,c.}
Wenn also zunächst $b$ oder $c = 0$ ist, so ist $a = d = 1$
und wir haben entweder
\[
A = \begin{pmatrix} 1, & b \\ 0, & 1 \end{pmatrix},
\]
was schon die Normalform $T$ hat, oder
\[
A = \begin{pmatrix} 1, & 0 \\ c, & 1 \end{pmatrix} = \begin{pmatrix} 0, & 1 \\ -1, & 0 \end{pmatrix} \begin{pmatrix} 1, & -c \\ 0, & 1 \end{pmatrix} \begin{pmatrix} 0, & -1 \\ 1, & 0 \end{pmatrix},
\]
so dass $\begin{pmatrix} 0, & -1 \\ 1, & 0 \end{pmatrix} A \begin{pmatrix} 0, & 1 \\ -1, & 0 \end{pmatrix}$ die Normalform $T$ hat. Sind aber $b$ und $c$ von Null verschieden,
so erhalten wir aus (8) die Transformation
\[
\begin{pmatrix} (d-1)\,b^{-1}, & 0, \\ -1, & (a-1)\,c^{-1} \end{pmatrix} \begin{pmatrix} a, & b \\ c, & d \end{pmatrix} \begin{pmatrix} (a-1)\,c^{-1}, & 0, \\ 1, & (d-1)\,b^{-1} \end{pmatrix} = \begin{pmatrix} 1, & -c \\ 0, & 1 \end{pmatrix},
\]
was die Form $T$ hat. Damit ist also bewiesen:

2.~\gesperrt{Eine Substitution $A$, in der $a+d = \pm 2$ ist, kann immer durch
Transformation auf die Normalform $T$ gebracht werden.}
